\begin{table}[H]
\centering
\scriptsize
\setlength{\tabcolsep}{2.6pt}
\begin{tabular}{lccc@{\hspace{8pt}}ccc}
\toprule
 & \multicolumn{3}{c}{CIFAR-100 images ($K{=}20$)} & \multicolumn{3}{c}{ImageNet images ($K{=}30$)} \\
\cmidrule(lr){2-4}\cmidrule(lr){5-7}
model & exc.\ ($z$) & depth ($z$) & $\hat\delta$ nat./Eucl. & exc.\ ($z$) & depth ($z$) & $\hat\delta$ nat./Eucl. \\
\midrule
MERU ViT-S & $-0.040$ (-3.0) & $+0.023$ (+1.7) & $0.163$/$0.163$ & $-0.024$ (-2.5) & $-0.004$ (-0.4) & $0.128$/$0.128$ \\
CLIP ViT-S & $-0.059$ (-4.9) & $-0.001$ (-0.1) & $0.140$/$0.141$ & $-0.025$ (-2.7) & $-0.005$ (-0.5) & $0.151$/$0.126$ \\
MERU ViT-B & $-0.042$ (-3.4) & $+0.023$ (+2.5) & $0.156$/$0.156$ & $-0.022$ (-2.4) & $+0.000$ (+0.0) & $0.129$/$0.129$ \\
CLIP ViT-B & $-0.056$ (-4.4) & $+0.010$ (+0.9) & $0.146$/$0.141$ & $-0.025$ (-2.3) & $-0.001$ (-0.1) & $0.151$/$0.124$ \\
MERU ViT-L & $-0.023$ (-1.5) & $+0.043$ (+3.1) & $0.170$/$0.171$ & $-0.021$ (-2.2) & $+0.004$ (+0.4) & $0.128$/$0.128$ \\
CLIP ViT-L & $-0.037$ (-2.9) & $+0.025$ (+2.1) & $0.170$/$0.155$ & $-0.025$ (-2.5) & $-0.003$ (-0.3) & $0.150$/$0.123$ \\
\bottomrule
\end{tabular}
\caption{Hyperbolic-backbone control. MERU \citep{desai2023meru} embeds images on the Lorentz hyperboloid with an entailment objective; its released Euclidean twin (a CLIP baseline: same ViT, same RedCaps data and recipe, minus the hyperbolic lift and entailment loss) differs only in geometry. Both pass through the paper's instrument unchanged (class centroids of the projected embeddings; ``exc.'' is the excess over the spectrum-matched null, ``depth'' the star-calibrated test of Table~\ref{tab:b29-depth}, ``nat.'' the Gromov $\hat\delta$ computed with the model's own metric: Lorentz distance between tangent-space-mean centroids for MERU, angular for CLIP). Training in hyperbolic space changes nothing at the class level: MERU shows the same clustering excess as its twin, no residual depth beyond a matched star on either dataset, and its Lorentz-metric $\hat\delta$ is indistinguishable from the Euclidean one. MERU's hierarchy is a generic$\to$specific (text$\supset$image) partial order, not a class taxonomy; the class-tree depth that no standard backbone shows is not present in a hyperbolic one either. The synthetic two-level tree of Table~\ref{tab:b25-star} remains the existence proof that the depth test fires when depth is present.}
\label{tab:b30-meru}
\end{table}
